% ===========================================================================
%  preamble.tex --- packages and macros for the pathogenesis chapter.
%  When this chapter is dropped into the thesis, merge the contents of this
%  file into the thesis preamble; duplicates may simply be deleted.
%
%  Derived from CH2/preamble.tex, the house preamble.  Where this chapter's
%  own main.tex disagreed with Chapter 2, Chapter 2 wins.
% ===========================================================================

\usepackage[T1]{fontenc}
\usepackage[utf8]{inputenc}
\usepackage{lmodern}
\usepackage{microtype}
\usepackage{amsmath,amssymb,amsthm,mathtools}
\usepackage{graphicx}
\usepackage{float}
\usepackage{placeins}
\usepackage{needspace}
\usepackage{booktabs}
\usepackage[hypcap=false,font=small,labelfont=bf]{caption}
\usepackage{subcaption}
\usepackage{tikz}
\usetikzlibrary{arrows.meta,positioning,calc}
\usepackage{pgfplots}
\pgfplotsset{compat=1.18}
\usepackage{enumitem}
\usepackage{xcolor}
\usepackage{array}
\usepackage{hyperref}
\usepackage[nameinlink,noabbrev]{cleveref}

\graphicspath{{figures/}}

% --- Equation cross-references as bare parenthesised numbers -----------------
\crefformat{equation}{(#2#1#3)}
\Crefformat{equation}{(#2#1#3)}
\crefrangeformat{equation}{(#3#1#4)--(#5#2#6)}
\crefmultiformat{equation}{(#2#1#3)}{ and (#2#1#3)}{, (#2#1#3)}{ and (#2#1#3)}

% --- Notation ---------------------------------------------------------------
% Scalars are set in italic; bold is reserved for vector-valued states, which
% appear only in the multi-type work of later chapters.  Vectors of unknowns
% in a linear system are the one exception, and are set upright bold.
%
% Symbols specific to this chapter, and the collisions that fixed them:
%   \varsigma   suppression rate per replicator-granule pair.  \alpha is the
%               Gompertz decay rate of the parked material and cell age in
%               the population chapter.
%   \vartheta   fixed removal count.  \theta is the exponential rate
%               \lambda(a-b) of the birth--death--catastrophe chapter.
%   \pi_{i,j}   extinction probability from (i,j).  \theta_{i,j} would
%               collide with \vartheta above.
%   Q_t         remaining granule count.  Y_t is the released count.
%   W_t         released count, matching the core chapters.
%   \mathcal{R} the per-cell replication ratio under removal.  R_0 is a
%               different object in the population chapter, where a theorem
%               proves it invariant under bursting.
%   \mathcal{M} truncated first moment.  \mathcal{K} is the burst size in the
%               three preceding chapters.
%   F(\vartheta) the burst-size c.d.f.  G is a generating function elsewhere.
% \delta (catastrophe rate) and \kappa are frozen thesis-wide and untouched.
\newcommand{\pgf}{probability generating function}
\newcommand{\rv}{random variable}
\newcommand{\pr}[1]{\mathbb{P}\!\left(#1\right)}
\newcommand{\ex}[1]{\mathbb{E}\!\left(#1\right)}
\newcommand{\var}[1]{\operatorname{Var}\!\left(#1\right)}
\newcommand{\dd}{\mathrm{d}}
\newcommand{\ee}{\mathrm{e}}
\newcommand{\ddt}[1]{\frac{\dd #1}{\dd t}}
\newcommand{\dda}[2]{\frac{\dd #1}{\dd #2}}
\newcommand{\pddt}[1]{\frac{\partial #1}{\partial t}}
\newcommand{\pdd}[2]{\frac{\partial #1}{\partial #2}}
\newcommand{\lrb}[1]{\left(#1\right)}
\newcommand{\lrbs}[1]{\left[#1\right]}
\newcommand{\ind}{\mathbf 1}
\newcommand{\ft}{\textit{F.~tularensis}}
\newcommand{\yp}{\textit{Y.~pestis}}

% --- Cross-chapter references -----------------------------------------------
% Redefining these macros is all that thesis integration requires: they become
% "Chapter~\ref{...}" once the chapter order is fixed.  No chapter number
% appears in body prose.
\newcommand{\ChM}{the mathematical introduction}
\newcommand{\ChCore}{the chapter on the birth--death--cata\-strophe process}
\newcommand{\ChCorecap}{The chapter on the birth--death--cata\-strophe process}
\newcommand{\ChDist}{the chapter on its distribution theory}
\newcommand{\ChDistcap}{The chapter on its distribution theory}
\newcommand{\ChPop}{the chapter on population dynamics}
\newcommand{\ChPopcap}{The chapter on population dynamics}
\newcommand{\ChTwoType}{the two-type chapter}
\newcommand{\ChEvo}{the chapter on variability in evolution}

% --- Forward references to later modelling chapters -------------------------
% Every forward reference is routed through \fwd, whose first argument is a
% stable key and whose second is the prose actually typeset.
\newcommand{\fwd}[2]{#2}

% --- Theorem environments ---------------------------------------------------
\theoremstyle{plain}
\newtheorem{theorem}{Theorem}[chapter]
\newtheorem{proposition}[theorem]{Proposition}
\newtheorem{lemma}[theorem]{Lemma}
\newtheorem{corollary}[theorem]{Corollary}
\theoremstyle{definition}
\newtheorem{definition}[theorem]{Definition}
\theoremstyle{remark}
\newtheorem{remark}[theorem]{Remark}

\allowdisplaybreaks

% --- Figure colours (house palette, matching Chapters 5 and 7) -------------
% Used by the state-space, fate-region and level-diagonal diagrams.
\definecolor{gray1}{HTML}{d9d9d9}
\definecolor{na1}{HTML}{0072B2}      % birth / cool
\definecolor{nb1}{HTML}{F0E442}      % contested (Okabe yellow)
\definecolor{mut1}{HTML}{D55E00}     % suppression / vermillion
\definecolor{Rec2}{HTML}{009E73}     % spent store / certain (teal)
\definecolor{Rup2}{HTML}{9A2820}     % extinction / catastrophe
\definecolor{blu}{HTML}{1A1C1F}      % ink
\definecolor{inksoft}{HTML}{565B62}
\definecolor{na1fill}{HTML}{E9F1F8}
\definecolor{Rec2fill}{HTML}{E6F4EE}
\definecolor{Rup2fill}{HTML}{F4E6E4}
